% !TEX root = task-6-4-3.tex
% Editable task bank. Append new BankTask blocks at the end of the document.
\documentclass[12pt,a4paper]{article}
\usepackage{physics-textbook}
\usepackage{graphicx,needspace}
\usetikzlibrary{decorations.pathmorphing}
\newif\ifShowSolutions
\ifdefined\TutorEdition
  \ShowSolutionstrue
\else
  \ShowSolutionsfalse % Change to \ShowSolutionstrue to display solutions.
\fi
\newcounter{banktask}
\newsavebox{\BankTaskBox}
\newcommand{\BankTaskContent}[3]{%
  \noindent\textbf{\thebanktask.}\quad #2\par
  \ifShowSolutions
    \Needspace{50mm}% Keep the solution heading with its opening text or diagram.
    #3\par
    \ifdefined\TutorEdition
      \ifcsname TutorSource#1\endcsname
        \nopagebreak\smallskip{\small\csname TutorSource#1\endcsname\par}%
      \fi
    \fi
  \fi
}
\newcommand{\BankTask}[3]{%
  \refstepcounter{banktask}\par\medskip
  \sbox{\BankTaskBox}{\begin{minipage}{\linewidth}%
    \BankTaskContent{#1}{#2}{#3}\end{minipage}}%
  \ifdim\dimexpr\ht\BankTaskBox+\dp\BankTaskBox\relax>\dimexpr\textheight-20mm\relax
    \Needspace{35mm}\BankTaskContent{#1}{#2}{#3}%
  \else
    \Needspace{\dimexpr\ht\BankTaskBox+\dp\BankTaskBox\relax}%
    \noindent\usebox{\BankTaskBox}\par
  \fi
  \medskip
}
\begin{document}
\BookSubsection{6.4.3}{Спектральный анализ}


% Source sdamgia-8870: https://phys-ege.sdamgia.ru/problem?id=8870
\BankTask{sdamgia-8870}{
На рисунке показаны спектры поглощения трех смесей неизвестных газов (Х, Y и Z), а также спектры излучения известных газов $1$ и $2$. Какая из смесей содержит газ $1$? 

\begin{center}\begin{tikzpicture}[book diagram,>=stealth]\draw(0,0.0)rectangle(5.6,0.6);\node[right]at(5.7,0.3){смесь газов X};\draw[line width=3.643pt](0.7171,0.03)--(0.7171,0.57);\draw[line width=1.943pt](1.3957,0.03)--(1.3957,0.57);\draw[line width=1.943pt](1.6006,0.03)--(1.6006,0.57);\draw[line width=1.943pt](3.1372,0.03)--(3.1372,0.57);\draw[line width=1.943pt](3.675,0.03)--(3.675,0.57);\draw(0,-0.8)rectangle(5.6,-0.20000000000000007);\node[right]at(5.7,-0.5){смесь газов Y};\draw[line width=1.943pt](0.3201,-0.77)--(0.3201,-0.2300000000000001);\draw[line width=1.943pt](0.525,-0.77)--(0.525,-0.2300000000000001);\draw[line width=1.943pt](1.4726,-0.77)--(1.4726,-0.2300000000000001);\draw[line width=1.943pt](1.6006,-0.77)--(1.6006,-0.2300000000000001);\draw[line width=2.672pt](4.1915,-0.77)--(4.1915,-0.2300000000000001);\draw[line width=2.915pt](4.9384,-0.77)--(4.9384,-0.2300000000000001);\draw[line width=2.429pt](5.0579,-0.77)--(5.0579,-0.2300000000000001);\draw(0,-1.6)rectangle(5.6,-1.0);\node[right]at(5.7,-1.3){смесь газов Z};\draw[line width=1.943pt](0.414,-1.57)--(0.414,-1.0300000000000002);\draw[line width=1.943pt](0.5165,-1.57)--(0.5165,-1.0300000000000002);\draw[line width=2.672pt](0.939,-1.57)--(0.939,-1.0300000000000002);\draw[line width=1.943pt](1.0543,-1.57)--(1.0543,-1.0300000000000002);\draw[line width=1.943pt](1.3872,-1.57)--(1.3872,-1.0300000000000002);\draw[line width=1.943pt](1.5921,-1.57)--(1.5921,-1.0300000000000002);\draw[line width=1.943pt](2.5909,-1.57)--(2.5909,-1.0300000000000002);\draw[line width=1.943pt](3.4104,-1.57)--(3.4104,-1.0300000000000002);\draw[line width=1.943pt](3.5128,-1.57)--(3.5128,-1.0300000000000002);\draw[line width=1.943pt](4.2043,-1.57)--(4.2043,-1.0300000000000002);\draw(0,-2.4000000000000004)rectangle(5.6,-1.8000000000000003);\node[right]at(5.7,-2.1000000000000005){газ 1};\draw[line width=2.672pt](0.939,-2.3700000000000006)--(0.939,-1.8300000000000005);\draw[line width=1.943pt](1.0543,-2.3700000000000006)--(1.0543,-1.8300000000000005);\draw[line width=2.429pt](2.5994,-2.3700000000000006)--(2.5994,-1.8300000000000005);\draw[line width=2.429pt](4.1957,-2.3700000000000006)--(4.1957,-1.8300000000000005);\draw(0,-3.2)rectangle(5.6,-2.6);\node[right]at(5.7,-2.9000000000000004){газ 2};\draw[line width=1.943pt](0.3116,-3.1700000000000004)--(0.3116,-2.6300000000000003);\draw[line width=2.186pt](0.5207,-3.1700000000000004)--(0.5207,-2.6300000000000003);\draw[line width=2.186pt](1.0073,-3.1700000000000004)--(1.0073,-2.6300000000000003);\draw[line width=2.186pt](1.3915,-3.1700000000000004)--(1.3915,-2.6300000000000003);\draw[line width=2.186pt](1.5963,-3.1700000000000004)--(1.5963,-2.6300000000000003);\draw[line width=2.186pt](3.1329,-3.1700000000000004)--(3.1329,-2.6300000000000003);\draw[line width=2.186pt](3.6707,-3.1700000000000004)--(3.6707,-2.6300000000000003);\end{tikzpicture}\end{center}
}{
\smallskip\noindent\textbf{Решение}\par
На рисунке представлены линейчатые спектры. Главное свойство линейчатого спектра состоит в том, линии какого-либо вещества на спектре зависят только от свойств атомов этого вещества, но совершенно не зависят от способа возбуждения свечения атомов. Из рисунка видно, что на спектре газовой смеси Z присутствуют линии, которые есть на спектре газа $1$. Этот факт говорит в пользу того, что смесь газа Z содержит газ $1$.

Ответ: Z.
}

\end{document}
